% =========================================================
% Proof: Intersection of Sigma-Algebras is a Sigma-Algebra
% Source: volume-i/algebras-of-sets/notes/...
% Mode: standard
% =========================================================

\subsection*{Intersection of $\sigma$-Algebras is a $\sigma$-Algebra}
\label{prf:intersection-sigma-algebras}

\begin{remark*}[Return]
\hyperref[prop:intersection-sigma-algebras]{$\leftarrow$ Back to Proposition (Intersection of $\sigma$-Algebras is a $\sigma$-Algebra) in Notes}
\end{remark*}

\begin{proof}
% TODO: proof to be filled.
% Strategy: Verify axioms ($\sigma$1)--($\sigma$3) pointwise across the intersecting family.
\end{proof}

\begin{remark*}[Proof shape]
% TODO

\FlashProof{Let $\mathcal{G} = \bigcap_{\alpha}\mathcal{F}_\alpha$. ($\sigma$1): $\emptyset\in\mathcal{F}_\alpha$ for all $\alpha$, so $\emptyset\in\mathcal{G}$. ($\sigma$2): if $A\in\mathcal{G}$ then $A\in\mathcal{F}_\alpha$ for all $\alpha$, so $A^c\in\mathcal{F}_\alpha$ for all $\alpha$, so $A^c\in\mathcal{G}$. ($\sigma$3): if each $A_n\in\mathcal{G}$ then each $A_n\in\mathcal{F}_\alpha$ for all $\alpha$, so $\bigcup A_n\in\mathcal{F}_\alpha$ for all $\alpha$, so $\bigcup A_n\in\mathcal{G}$.}
\end{remark*}

\begin{remark*}[Dependencies]
\begin{itemize}
  \item \hyperref[def:sigma-algebra]{Definition of $\sigma$-algebra}: axioms ($\sigma$1)--($\sigma$3)
\end{itemize}

\FlashUses{def:sigma-algebra}
\end{remark*}
